\documentclass[10pt,a4paper]{article}

\usepackage{fontspec}
\usepackage[margin=2cm]{geometry}
\usepackage{xcolor}
\usepackage{array,booktabs,tabularx,colortbl}
\usepackage{amsmath,amssymb}
\usepackage{enumitem}
\usepackage{fancyhdr}
\usepackage{titlesec}
\usepackage[most]{tcolorbox}
\usepackage{parskip}
\usepackage{hyperref}

% ── Paleta ──
\definecolor{azul}{HTML}{185FA5}
\definecolor{azulclr}{HTML}{E6F1FB}
\definecolor{teal}{HTML}{0F6E56}
\definecolor{tealclr}{HTML}{E1F5EE}
\definecolor{gris}{HTML}{5F5E5A}
\definecolor{grisclr}{HTML}{F1EFE8}
\definecolor{coral}{HTML}{993C1D}
\definecolor{negro}{HTML}{2C2C2A}
\definecolor{alerta}{HTML}{D32F2F}
\definecolor{alertaclr}{HTML}{FFEBEE}

% ── Fuente sans-serif ──
\setmainfont{Arial}
\setsansfont{Arial}
\setmonofont{Liberation Mono}

\renewcommand{\familydefault}{\sfdefault}

% ── Estilos ──
\titleformat{\section}{\bfseries\Large\color{azul}}{}{0em}{}[\vspace{-2pt}]
\titleformat{\subsection}{\bfseries\normalsize\color{teal}}{}{0em}{}

\pagestyle{fancy}
\fancyhf{}
\fancyhead[C]{\color{gris}\small\textbf{Hoja de Fórmulas de Cinemática — Unidad 2}}
\fancyfoot[C]{\color{gris}\thepage}
\renewcommand{\headrule}{\color{azul}\rule{\textwidth}{0.4pt}}

% ── Macros ──
\newcommand{\formulabox}[2]{%
  \begin{tcolorbox}[colback=azulclr, colframe=azul, arc=4pt, boxrule=0.5pt,
    left=8pt, right=8pt, top=6pt, bottom=6pt, width=\textwidth]
    \textbf{\textcolor{azul}{#1}} \quad $\displaystyle #2$
  \end{tcolorbox}%
}

\newcommand{\alertbox}[1]{%
  \begin{tcolorbox}[colback=alertaclr, colframe=alerta, arc=4pt, boxrule=0.5pt,
    left=8pt, right=8pt, top=6pt, bottom=6pt, width=\textwidth]
    \textbf{\textcolor{alerta}{#1}}
  \end{tcolorbox}%
}

\newcommand{\infobox}[1]{%
  \begin{tcolorbox}[colback=grisclr, colframe=gris, arc=4pt, boxrule=0.3pt,
    left=8pt, right=8pt, top=4pt, bottom=4pt, width=\textwidth]
    \small\textcolor{gris}{#1}
  \end{tcolorbox}%
}

\newcolumntype{L}{>{\raggedright\arraybackslash}X}
\newcolumntype{C}{>{\centering\arraybackslash}X}

\begin{document}

% ═════════════════════════════════════════════════
% ENCABEZADO
% ═════════════════════════════════════════════════
\thispagestyle{fancy}
\begin{center}
  {\Huge\bfseries\color{azul} HOJA DE FÓRMULAS DE CINEMÁTICA}\\[4pt]
  {\large\color{gris} Unidad 2: Correcciones y Complementos}\\[4pt]
  {\color{gris} $g = 9.8\ \text{m/s}^2$ \;|\; Sistema Internacional de Unidades (SI)}
\end{center}
\vspace{6pt}
\hrule height 0.5pt \color{gris}
\vspace{10pt}

% ═════════════════════════════════════════════════
% CORRECCIÓN: DISTANCIA VS DESPLAZAMIENTO
% ═════════════════════════════════════════════════
\alertbox{\textbf{CORRECCIÓN IMPORTANTE:} Distancia vs.\ Desplazamiento}

En los apuntes originales se definía la velocidad como ``desplazamiento / tiempo'', pero se usaba $d = \text{distancia}$ en la fórmula. Esto es un error conceptual. La fórmula correcta es:

\formulabox{Velocidad (vectorial)}{v = \frac{\Delta d}{t} \quad (\text{usa DESPLAZAMIENTO, cambio neto de posición})}
\formulabox{Rapidez (escalar)}{v = \frac{d}{t} \quad (\text{usa DISTANCIA, espacio total recorrido})}

\vspace{6pt}
\hrule height 0.5pt
\vspace{12pt}

% ═════════════════════════════════════════════════
% 1. MRU
% ═════════════════════════════════════════════════
\section{1. Movimiento Rectilíneo Uniforme (MRU)}

La velocidad es constante, por lo tanto, la aceleración es cero ($a = 0$).

\vspace{4pt}
\begin{tabularx}{\textwidth}{LCC}
  \toprule
  \rowcolor{azul}
  \textbf{\textcolor{white}{Magnitud}} &
  \textbf{\textcolor{white}{Fórmula}} &
  \textbf{\textcolor{white}{Cuándo usarla}} \\
  \midrule
  Velocidad     & $v = \Delta d / t$     & Conocemos desplazamiento y tiempo \\
  \rowcolor{azulclr}
  Desplazamiento & $\Delta d = v \cdot t$ & Conocemos velocidad y tiempo \\
  Tiempo        & $t = \Delta d / v$     & Conocemos desplazamiento y velocidad \\
  \bottomrule
\end{tabularx}
\vspace{12pt}

% ═════════════════════════════════════════════════
% 2. MRUA
% ═════════════════════════════════════════════════
\section{2. Movimiento Rectilíneo Uniformemente Acelerado (MRUA)}

La aceleración es constante. Estas son las 4 ecuaciones fundamentales que faltaban en los apuntes:

\vspace{4pt}
\begin{tabularx}{\textwidth}{CCC}
  \toprule
  \rowcolor{azul}
  \textbf{\textcolor{white}{Ecuación}} &
  \textbf{\textcolor{white}{Fórmula}} &
  \textbf{\textcolor{white}{Falta conocer...}} \\
  \midrule
  1ª Ecuación & $V_f = V_i + a \cdot t$          & El desplazamiento ($d$) \\
  \rowcolor{azulclr}
  2ª Ecuación & $d = V_i \cdot t + \frac{1}{2} a \cdot t^2$ & La velocidad final ($V_f$) \\
  3ª Ecuación & $V_f^2 = V_i^2 + 2 a \cdot d$   & El tiempo ($t$) \\
  \rowcolor{azulclr}
  4ª Ecuación & $d = \dfrac{V_i + V_f}{2} \cdot t$ & La aceleración ($a$) \\
  \bottomrule
\end{tabularx}
\vspace{12pt}

% ═════════════════════════════════════════════════
% 3. CAÍDA LIBRE
% ═════════════════════════════════════════════════
\section{3. Caída Libre}

Movimiento bajo la única influencia de la gravedad ($g = 9.8\ \text{m/s}^2$ hacia abajo). Fórmulas para un objeto que se deja caer desde el reposo ($V_i = 0$):

\vspace{4pt}
\begin{tabularx}{\textwidth}{LC}
  \toprule
  \rowcolor{azul}
  \textbf{\textcolor{white}{Magnitud}} &
  \textbf{\textcolor{white}{Fórmula}} \\
  \midrule
  Velocidad final   & $V_f = g \cdot t$ \\
  \rowcolor{azulclr}
  Altura caída      & $h = \dfrac{1}{2} \cdot g \cdot t^2$ \\
  Velocidad vs Altura & $V_f^2 = 2 \cdot g \cdot h$ \\
  \rowcolor{azulclr}
  Tiempo de caída   & $t = \sqrt{\dfrac{2h}{g}}$ \\
  \bottomrule
\end{tabularx}
\vspace{12pt}

% ═════════════════════════════════════════════════
% 4. TIRO VERTICAL
% ═════════════════════════════════════════════════
\section{4. Tiro Vertical}

Cuando se lanza un objeto hacia arriba, la gravedad actúa en sentido contrario al movimiento (desaceleración). En el punto más alto, la velocidad es momentáneamente cero ($V_f = 0$).

\vspace{4pt}
\begin{tabularx}{\textwidth}{LL}
  \toprule
  \rowcolor{teal}
  \textbf{\textcolor{white}{Fase}} &
  \textbf{\textcolor{white}{Fórmulas}} \\
  \midrule
  Subida (desacelera) &
  $V_f = V_i - g \cdot t \quad\mid\quad
   h = V_i \cdot t - \frac{1}{2} g \cdot t^2 \quad\mid\quad
   V_f^2 = V_i^2 - 2 g \cdot h$ \\
  \rowcolor{tealclr}
  Punto más alto &
  $V_f = 0\ \text{m/s} \quad\mid\quad
   t_{\text{subida}} = \dfrac{V_i}{g} \quad\mid\quad
   H_{\text{max}} = \dfrac{V_i^2}{2g}$ \\
  Bajada (acelera) &
  $V_f = g \cdot t_{\text{bajada}} \quad\mid\quad
   h = \dfrac{1}{2} \cdot g \cdot t_{\text{bajada}}^2$ \\
  \bottomrule
\end{tabularx}
\vspace{12pt}

% ═════════════════════════════════════════════════
% 5. TIRO PARABÓLICO
% ═════════════════════════════════════════════════
\section{5. Tiro Parabólico}

Es la superposición de un MRU en el eje X y un MRUA en el eje Y. Lo primero que se debe hacer es descomponer la velocidad inicial, algo que faltaba en los apuntes originales:

\vspace{4pt}
\begin{tabularx}{\textwidth}{LL}
  \toprule
  \rowcolor{azul}
  \textbf{\textcolor{white}{Componente}} &
  \textbf{\textcolor{white}{Fórmula}} \\
  \midrule
  Velocidad inicial X & $V_{0x} = V_0 \cdot \cos(\theta)$ \hfill (Se mantiene constante) \\
  \rowcolor{azulclr}
  Velocidad inicial Y & $V_{0y} = V_0 \cdot \sin(\theta)$ \hfill (Cambia por la gravedad) \\
  Posición en X (MRU) & $x = V_{0x} \cdot t$ \\
  \rowcolor{azulclr}
  Posición en Y (MRUA) & $y = V_{0y} \cdot t - \dfrac{1}{2} \cdot g \cdot t^2$ \\
  Velocidad en Y & $V_y = V_{0y} - g \cdot t$ \\
  \rowcolor{azulclr}
  Tiempo de vuelo total & $T = \dfrac{2 \cdot V_0 \cdot \sin(\theta)}{g}$ \\
  Altura máxima & $H = \dfrac{V_0^2 \cdot \sin^2(\theta)}{2g}$ \\
  \rowcolor{azulclr}
  Alcance horizontal & $R = \dfrac{V_0^2 \cdot \sin(2\theta)}{g}$ \\
  \bottomrule
\end{tabularx}
\vspace{12pt}

% ═════════════════════════════════════════════════
% RECORDATORIOS DE CONVERSIÓN
% ═════════════════════════════════════════════════
\vspace{6pt}
\hrule height 0.5pt
\vspace{12pt}

\section*{Recordatorios de Conversión de Unidades}
\begin{itemize}[leftmargin=*, nosep]
  \item De km/h a m/s: dividir entre 3.6 \hfill (ej: $72\ \text{km/h} \div 3.6 = 20\ \text{m/s}$)
  \item De m/s a km/h: multiplicar por 3.6 \hfill (ej: $15\ \text{m/s} \times 3.6 = 54\ \text{km/h}$)
  \item De minutos a segundos: multiplicar por 60
\end{itemize}

\vspace{8pt}
\infobox{\textbf{Nota:} Esta hoja complementa los apuntes de la Unidad 2, corrigiendo las omisiones y errores conceptuales para garantizar que los cálculos en los ejercicios se realicen con las ecuaciones físicas correctas.}

\end{document}
